% 移除了导致错误的 magic comment
\documentclass[UTF8,a4paper,zihao=5,fontset=fandol]{ctexart}

% 本模板依据《经济研究》网站 2025-01-03 发布的“稿件体例要求”整理。
% 网页已明确的字体字号要求已尽量固化；页边距和行距网页未明示，以下采用中性设置，可按编辑部反馈微调。
% 建议使用 XeLaTeX 编译。

\usepackage[a4paper,left=2.8cm,right=2.8cm,top=2.8cm,bottom=2.6cm,footskip=1.2cm]{geometry}
\usepackage{fontspec}
\usepackage{amsmath,amssymb,bm}
\usepackage{array,booktabs,longtable,tabularx,threeparttable}
\usepackage{graphicx}
\usepackage{url}
\usepackage{caption}
\usepackage{setspace}
\usepackage{indentfirst}
\usepackage{titlesec}
\usepackage{perpage}

\setmainfont{Times New Roman}
\defaultfontfeatures{Ligatures=TeX}

\MakePerPage{footnote}

\setlength{\parindent}{2em}
\setlength{\parskip}{0pt}
\setstretch{1.45}
\pagestyle{plain}

\makeatletter
\renewcommand\footnotesize{\songti\zihao{6}}
\renewcommand{\footnoterule}{%
  \vspace{0.6\baselineskip}%
  \hrule width 0.18\textwidth
  \vspace{0.3\baselineskip}%
}
\makeatother

\newcommand{\ERJTitleSize}{\zihao{-3}} % 需要三号时改为 \zihao{3}
\newcommand{\ERJAuthorSize}{\zihao{-4}}

\DeclareCaptionFont{erjfigurefont}{\heiti\zihao{-5}}
\DeclareCaptionFont{erjtablefont}{\songti\zihao{5}}
\captionsetup{justification=centering,singlelinecheck=false}
\captionsetup[figure]{font=erjfigurefont,labelsep=space}
\captionsetup[table]{font=erjtablefont,labelsep=space,position=top}

\newcolumntype{L}[1]{>{\raggedright\arraybackslash}p{#1}}
\newcolumntype{C}[1]{>{\centering\arraybackslash}p{#1}}
\newcommand{\ERJTableFont}{%
  \fangsong\zihao{-5}%
  \setlength{\tabcolsep}{2pt}%
  \renewcommand{\arraystretch}{1.16}%
}
\newcommand{\ERJTableNoteFont}{\songti\zihao{6}}

\renewcommand{\thesection}{\chinese{section}}
\renewcommand{\thesubsection}{（\chinese{subsection}）}
\renewcommand{\thesubsubsection}{\arabic{subsubsection}.}
\renewcommand{\theparagraph}{(\arabic{paragraph})}
\setcounter{secnumdepth}{4}

\titleformat{\section}
  {\centering\songti\zihao{-4}}
  {\thesection、}
  {0.5em}
  {}
\titlespacing*{\section}{0pt}{1.4\baselineskip}{0.8\baselineskip}

\titleformat{\subsection}[hang]
  {\songti\zihao{5}}
  {\hspace*{2em}\thesubsection}
  {0.3em}
  {}
\titlespacing*{\subsection}{0pt}{0.9\baselineskip}{0.3\baselineskip}

\titleformat{\subsubsection}[hang]
  {\songti\zihao{5}}
  {\hspace*{2em}\thesubsubsection}
  {0.3em}
  {}
\titlespacing*{\subsubsection}{0pt}{0.7\baselineskip}{0.2\baselineskip}

\titleformat{\paragraph}[runin]
  {\songti\zihao{5}}
  {\hspace*{2em}\theparagraph}
  {0.3em}
  {}
\titlespacing*{\paragraph}{0pt}{0.5\baselineskip}{0.5em}

\newcommand{\erjthanks}[1]{\thanks{\hspace*{2em}#1}}
\newcommand{\erjfootnote}[1]{\footnote{\hspace*{2em}#1}}

\newcommand{\ERJAbstract}[1]{%
  \par\noindent{\fangsong\zihao{5}内容摘要：#1\par}%
}

\newcommand{\ERJKeywords}[1]{%
  \par\noindent{\fangsong\zihao{5}关键词：#1\par}%
}

\newcommand{\ERJWhere}[1]{%
  \par\noindent\hspace*{2em}{\songti\zihao{5}其中，#1\par}%
}

\newcommand{\coefse}[2]{%
  \begin{tabular}{@{}c@{}}#1\\[-0.1em] (#2)\end{tabular}%
}

\newcommand{\ERJRef}[1]{%
  \par\noindent\hspace*{2em}\hangindent=2em\hangafter=1 {#1}\par%
}

\makeatletter
\renewcommand{\maketitle}{%
  \begingroup
  \renewcommand{\thefootnote}{\fnsymbol{footnote}}%
  \begin{center}
    {\songti\ERJTitleSize \@title\par}
    \vspace{1.1em}
    {\songti\ERJAuthorSize \@author\par}
  \end{center}
  \vspace{1.0em}
  \@thanks
  \endgroup
  \setcounter{footnote}{0}%
  \renewcommand{\thefootnote}{\arabic{footnote}}%
}
\makeatother

\title{地方债务压力如何影响政府产业引导基金的创新效应？——来自地级市层面的证据}
\author{%
  \renewcommand{\arraystretch}{1.18}%
  \begin{tabular}{@{}c@{\hspace{1.8em}}c@{\hspace{1.8em}}c@{\hspace{1.8em}}c@{\hspace{1.8em}}c@{}}
    周子川 & 张晋轩 & 余阳 & 张静妍 & 陈芊羽 \\
    {\zihao{5}2411597} & {\zihao{5}2414045} & {\zihao{5}2412782} & {\zihao{5}2411588} & {\zihao{5}2412405}
  \end{tabular}
}
\date{}

\begin{document}

\maketitle

\ERJAbstract{政府产业引导基金能否稳定促进城市创新，以及地方债务压力会否削弱其政策绩效，是理解政府投资基金高质量发展的重要问题。本文基于2015—2024年地级市—年份面板数据，采用双向固定效应模型，考察政府产业引导基金对城市创新的影响及地方债务压力的约束作用。研究发现：政府产业引导基金总体上显著促进城市创新，其中累计设立规模这一更能表征持续资本供给能力的口径作用更为明显；进一步看，地方债务压力显著削弱了基金的创新促进效应，该结论在发明申请、实用新型申请和专利申请总量等口径下均成立。机制分析表明，债务压力主要通过压缩基金投向早期项目的空间、削弱社会资本撬动效率以及弱化基金改善地区金融发展环境的能力，进而抑制其创新绩效。本文表明，政府产业引导基金的创新效应并非无条件成立，其政策绩效在很大程度上取决于地方债务约束和区域金融环境，因而推动政府投资基金高质量发展，应统筹基金政策性定位、债务风险防控与科技金融体系建设。}
\ERJKeywords{政府产业引导基金；债务压力；城市创新；耐心资本；社会资本撬动效率；融资约束}


% ===== BEGIN INPUT: 一_引言.tex =====
% 本文件由 main.tex 输入，作为“引言”章节，不单独编译。

\section{引言}

创新是推动高质量发展的重要动力，科技创新能力提升已经成为现代化产业体系建设中的关键内容。近几年，围绕科技创新、产业升级和新质生产力培育，政府投资基金和创业投资在政策层面受到持续重视。《国务院办公厅关于印发〈促进创业投资高质量发展的若干政策措施〉的通知》提出，要围绕创业投资“募投管退”全链条完善政策环境，强化企业创新主体地位，引导创业投资更好发挥投早、投小、投长期、投硬科技的作用。随后，《国务院办公厅关于促进政府投资基金高质量发展的指导意见》进一步明确，政府投资基金要突出政府引导和政策性定位，按照市场化、法治化、专业化原则规范运作，发展耐心资本，更好服务现代化产业体系建设，同时还强调要合理统筹基金布局，防止同质化竞争和对社会资本产生挤出效应。在这样的政策背景下，政府产业引导基金已经不只是一般意义上的财政工具，而是被赋予了支持科技创新、撬动社会资本和优化产业结构的重要功能。

从现实发展来看，政府产业引导基金扩张较快，已经成为地方政府支持创新和招商引资的重要手段。已有研究表明，政府引导基金可以通过提高失败容忍度、缓解融资约束和集聚社会资本等方式促进企业创新，尤其是在关键核心技术突破、区域创新投入增长等方面表现出一定的积极作用（吴超鹏和严泽浩，2023；蔡庆丰等，2024）。不过，政府产业引导基金的政策目标能否稳定实现，并不是一个没有条件的问题。与政策文本中对“耐心资本”“投早投小”和“市场化运作”的强调相比，地方政府在现实中还面临着财政收支矛盾加大、债务风险防控趋严和预算约束增强等压力。特别是《国务院办公厅关于促进政府投资基金高质量发展的指导意见》已经明确提出，政府投资基金布局和运行要统筹考虑财力基础、产业基础以及债务风险，这说明地方债务约束并不是外在背景，而是已经进入政府投资基金高质量发展的政策框架之中。

也正因为如此，一个值得进一步讨论的问题就变得比较清楚了：当地方政府同时承担稳增长、促创新和防风险等多重任务时，债务压力会不会改变政府产业引导基金的运作逻辑，进而影响其对创新的促进作用。按政策设计看，政府产业引导基金应当通过股权投资方式支持高风险、长周期的创新活动，并通过让利和增信机制撬动更多社会资本进入科技创新领域；但在债务压力较大的情形下，地方政府可能更重视基金安全性、回收速度和短期绩效，这样一来，基金原本应有的耐心资本属性、社会资本撬动功能以及融资约束缓解作用，就有可能被削弱。换句话说，政府产业引导基金是否促进创新，不能只从基金本身来判断，还需要把它放到地方财政债务约束不断增强的现实情境中来考察。

现有研究已经分别从两个方面积累了较多成果。一方面，关于政府产业引导基金与创新的研究表明，引导基金既可能通过融资支持、信号传递和资本集聚等机制促进创新，也可能因为代理问题、风险规避和目标偏离而削弱创新激励。另一方面，关于财政压力和地方政府债务的研究则指出，财政收支约束和债务扩张会通过挤出金融资源、提高融资成本、强化短期激励和扭曲政府干预方式等渠道抑制企业创新与长期投资。相比之下，把“地方债务压力—政府产业引导基金运作—创新绩效”放在同一分析框架中的研究还比较少。现有文献更多是在讨论基金是否有效，或者财政压力是否抑制创新，而对债务压力如何改变政府基金的行为边界，特别是如何影响其耐心资本属性、社会资本撬动效率和融资约束缓解功能，仍缺少较为系统的分析。

基于此，本文将地方债务压力、政府产业引导基金与创新绩效纳入统一分析框架，重点讨论两个层面的问题：第一，政府产业引导基金总体上能否促进创新；第二，在地方债务压力上升的背景下，这种促进作用是否会被削弱，以及这种削弱主要通过哪些机制实现。围绕上述问题，本文进一步从耐心资本属性、社会资本撬动效率和融资约束缓解效应三个方面展开分析，力图更完整地说明政府产业引导基金为什么在某些地区和条件下能够较好促进创新，而在另一些地区和条件下其效果又会明显减弱。

本文的研究可能有以下几方面意义。第一，现有研究多把政府产业引导基金视为独立的创新政策工具，本文则进一步把地方债务压力纳入分析，有助于从财政约束角度理解基金创新效应的差异来源。第二，已有研究虽然讨论了财政压力抑制创新的问题，但对政府产业引导基金这一中观政策工具的传导作用关注不多，本文有助于拓展地方债务压力影响创新的机制链条。第三，从现实政策角度看，两份国务院文件一方面强调发展耐心资本、支持科技创新，另一方面也强调统筹债务风险、防止同质化竞争和挤出效应，本文的讨论能够在一定程度上回应这种政策上的内在张力，并为理解政府产业引导基金何以“有效”、又何以可能“失效”提供补充证据。

下文其余部分安排如下：第二部分为文献综述与理论假说，第三部分为研究设计，第四部分为实证结果分析，第五部分为结论及进一步讨论。
% ===== END INPUT: 一_引言.tex =====



% ===== BEGIN INPUT: 二_文献综述与理论假说.tex =====
% 本文件由 main.tex 输入，作为“文献综述与理论假说”章节，不单独编译。

\section{文献综述与理论假说}

\subsection{文献综述}

\subsubsection{基金促进创新的效应与机制的已有研究}

围绕政府产业引导基金能否促进创新，现有研究大体形成了“促进论”与“约束论”两类观点。较多研究认为，引导基金本质上是财政资金“基金化”和产业政策“金融化”的重要载体，其通过股权投资而非传统补贴方式介入企业创新活动，有助于克服科技创新中的融资约束、信息不对称和风险外部性，从而提升企业创新投入与创新产出。吴超鹏和严泽浩（2023）发现，政府基金引导能够显著促进企业在关键核心技术领域的创新，其重要原因在于政府资本提高了对创新失败的容忍度，并缓解了企业外部融资约束。刘玉斌等（2023）基于战略性新兴产业创业投资引导基金的研究也表明，政府引导基金对企业创新具有显著促进作用，其机制主要体现为融资约束缓解、认证效应与信号传递。蔡庆丰等（2024）则进一步从区域层面指出，产业引导基金不仅能够直接促进域内企业创新投入，而且能够通过资本集聚形成融资效应、信息效应和竞争效应，放大财政资金“拨改投”的微观创新激励。

不过，另一部分研究对引导基金的创新绩效持更为审慎的判断。这类研究强调，政府资本虽然具有政策目标，但也可能因为多重委托代理、风险规避、考核约束和寻租行为而偏离“投早、投小、投科技”的初衷。徐明（2022）发现，政府风险投资不仅未必促进企业创新，反而可能因目标冲突、治理失灵和政治关联而抑制企业研发投入和创新产出。进一步地，部分文献指出，引导基金在实际运作中可能偏好相对成熟、风险较低的项目，由此形成对民间风险资本的挤出或同质化竞争，而非真正填补创新融资缺口。由此可见，现有研究并不否认引导基金具有促进创新的潜在能力，但其创新绩效能否实现，很大程度上取决于基金是否能够保持耐心资本属性、有效撬动社会资本，并真正流向高风险、长周期的创新活动。

总体来看，关于“基金是否促进创新”的争论，已经从简单的政策有效性判断，转向对其作用条件和内在机制的识别。这为本文进一步讨论债务约束条件下政府产业引导基金何以促进创新、又何以可能偏离创新目标，提供了直接的文献基础。

\subsubsection{地方财政债务压力对创新的已有研究}

与引导基金研究相对应，另一条重要文献线索聚焦于地方财政压力和政府债务扩张对创新活动的影响。既有研究普遍认为，财政压力与债务约束并非单纯的宏观背景变量，而是会通过改变政府行为、金融资源配置和企业融资环境，对创新形成系统性影响。李森和王聪（2024）发现，财政压力会显著抑制制造业企业创新，并通过加剧政治晋升激励、减少创新补贴和提升企业税负等渠道发挥作用。类似地，部分研究指出，在财政收支矛盾加剧时，地方政府更可能强化税收征管、压缩长期性科技支持支出，从而弱化企业开展高风险研发活动的能力和意愿。

从债务扩张的角度看，现有研究更强调其金融挤出效应与资源错配后果。余海跃和康书隆（2020）发现，地方政府债务扩张会通过抬升企业融资成本挤出企业投资，且对高效率民营企业的影响更强。徐彦坤（2020）的研究进一步表明，地方政府债务上升会加剧企业融资约束、抬高债务资本成本并压缩投资机会。刘贯春等（2022）则从投融资期限结构角度指出，地方政府债务治理有助于缓解政企融资竞争并改善企业“短贷长投”问题，其反向含义是，债务扩张本身会恶化长期资金供给。除融资挤出外，地方政府债务还可能改变企业的资产配置方向。孔丹凤和谢国梁（2021）发现，地方政府债务扩张会加剧企业金融化倾向，推动资金从实体创新活动转向短期金融收益。

更进一步，财政压力和债务压力还会改变地方政府干预市场的方式。张跃文等（2025）指出，在财政压力上升背景下，地方政府可能通过“合作型”干预影响企业投资决策，以增加短期税源和财政收入。这说明，财政压力不仅影响政府“有没有钱支持创新”，还会影响政府“如何使用政策工具”。如果把这一逻辑延伸到政府产业引导基金，就意味着地方债务压力可能促使地方政府更加重视资金安全、回收速度和短期绩效，进而削弱基金原本应有的耐心资本属性、让利属性和创新导向属性。

因此，现有关于财政债务压力的研究已经揭示出一个重要事实：地方政府面临的债务约束会同时影响企业创新环境和政府自身的行为逻辑。只是现有文献更多考察财政压力或地方债务对企业创新的直接影响，对于其如何通过扭曲政府产业引导基金这一政策工具而间接影响创新，仍缺乏专门讨论。

\subsubsection{现有文献的不足与本文边际贡献}

综合来看，现有研究虽然分别从政府产业引导基金和地方财政债务压力两个维度积累了较为丰富的成果，但仍存在三方面不足。第一，两条文献线索大多彼此分离。已有研究要么关注引导基金能否促进企业创新，要么考察财政压力和地方债务如何抑制企业创新，较少将“地方债务压力—引导基金运作—企业创新绩效”纳入统一分析框架，从而难以解释为何同样是政府产业引导基金，在不同财政约束情境下可能表现出显著差异。第二，既有研究对于引导基金失灵的讨论，虽然涉及风险规避、代理问题和资本挤出，但较少进一步从耐心资本属性、社会资本撬动效率和融资约束缓解效应三条机制链进行系统展开，因而对债务压力如何改变基金行为逻辑的解释仍不够细化。第三，现有研究更多停留在“基金是否有效”的平均效应判断上，对政策工具的作用边界讨论不足，尤其缺少对地方债务压力这一现实制度约束的关注。

基于上述不足，本文尝试将地方财政债务压力、政府产业引导基金与创新绩效纳入统一分析框架，重点考察债务压力如何削弱政府产业引导基金促进创新的作用，并从耐心资本属性扭曲、社会资本撬动效率下降和融资约束缓解效应弱化三个方面展开理论分析。相较于既有文献，本文的边际贡献主要体现在：其一，不再停留于讨论“引导基金是否促进创新”的一般性命题，而是进一步回答在债务约束背景下这一促进作用为何会减弱；其二，将宏观债务约束与中观政策工具运作结合起来，拓展了地方债务压力影响创新的传导链条研究；其三，从政策工具作用边界的角度，为理解政府产业引导基金何以“有效”、又何以“失效”提供新的分析视角。基于此，下文将进一步提出本文的理论假说。

\subsection{理论分析和理论假说}

企业的创新活动因其长期性和不确定性，受到严重的融资约束，主要依靠企业内部资金( Hall，2002)。但企业利用内部资金进行创新活动存在两大难题：一是内部财务波动会直接影响创新投入，易受到外部冲击，可能导致资金链断裂，创新活动被迫中断；二是创新活动具有较高的调整成本，突然中断和再重新启动，会使企业承担巨大损失（Hall，2002，2005）。为解决这一问题，受融资约束的企业会根据调整成本的高低进行反向调整，优先削减调整成本低的投资，将有限资金保留用于调整成本高的创新活动，最后保持各种投资的净收益在边际上相等，从而保证创新投资的连续性与稳定性( Fazzari et al.1993)。

然而，初创期和早中期的科技型企业普遍存在企业创新资金不足、外部融资成本过高、融资渠道不通畅的问题。从市场失灵理论来看，战略性新兴产业企业的创新活动面临三重典型的市场失灵：一是信息不对称，企业技术前景不确定性高、投资周期长的特点，外部投资者需付出较高成本以准确评估其偿债能力与成长潜力；二是正外部性，企业创新成果存在知识溢出效应，社会收益高于私人收益，市场化资本在收益风险权衡下天然倾向于投资不足；三是融资缺口，企业初创期和早中期缺乏稳定的现金流与可抵押资产。在这一背景下，企业内部可以用来平滑创新波动的资本十分有限，单靠市场机制作用与企业自身管理难以有效缓解其面临的较强融资约束。

为应对这一问题，政府产业引导基金作为一项纠正市场失灵的政策性金融安排，为战略性新兴产业企业突破融资约束提供了可行路径。政府产业引导基金参股基金根据企业所处生命周期的不同阶段，实施阶段性投资与联合投资，匹配企业差异化的资金需求，直接为企业创新活动提供资本支持。同时，政府产业引导基金的发展有利于促进地区金融资源的均衡分布，资本资源聚集的网络效应有利于企业信息在资本之间的共享，降低社会资本寻求投资对象过程中的信息成本（李健旋和赵林度，2018）。通过上述机制，政府产业引导基金不仅直接补充了企业创新资金，更改善了企业融资条件，为社会资本进入创新领域提供保障，同时有助于企业保持创新投资的稳定性与连续性。

由此提出假说H1：政府产业引导基金能够促进城市创新。

然而，引导基金的杠杆运作模式在放大财政资金效果的同时，近年来也面临地方债务压力日益严峻的现实问题。截至2024年末，地方政府法定债务余额已达47.54万亿元，隐性债务约14.3万亿元，地方整体偿债压力持续攀升（国务院，2025）。Wang等（2025）在《Sustainability》上的研究发现，地方政府债务扩张显著降低企业创新质量，且财政压力较低的地区这一负面影响更小。Zhang和Jin（2022）发现地方债务与创新的倒U型关系，拐点为债务与GDP之比0.689。这些发现告诉我们，债务压力可能改变了引导基金支持创新产出的制度环境和运作条件。债务压力可能会通过恶化融资约束，侵蚀耐心资本属性，减少社会资本撬动作用等途径，影响产业引导基金支持创新的效果和能力。

由此提出假说H2:地方债务压力会削弱政府产业引导基金的创新促债务压力调节效应。

创新活动的长周期性、高不确定性和正外部性这三个特征共同决定了创新活动对耐心资本的内生需求：一种能够长期持有、耐受不确定性、不追求短期回报退出的资本形态。理解耐心资本作为引导基金的本质，需要区分两个层次。第一个层次是功能层次：引导基金填补市场耐心资本的供给缺口，在市场资本“不愿投、不敢投”的领域发挥补位作用。第二个层次是信号层次：引导基金的存在本身就是对创新长期价值的制度性背书，其“耐心”属性向市场传递了一个关键信号——某些创新项目的价值需要以十年为单位衡量，而非以季度或年度衡量。这个信号具有公共品性质，私人资本不会自发提供。

引导基金对地方创新产出的支持效果，从根本上取决于其耐心资本属性能否完整保持, 将资金投向国务院要求的“投早投小投科技”的方向。一旦这一属性被侵蚀，引导基金就不再能够履行其制度功能，其支持创新产出的效果必然下降。

首先，偿债支出具有预算优先权，在财政资源有限时优先于引导基金出资。这可能会导致新设基金被推迟或取消，已有基金的后续出资被延缓，并且基金存续期被迫缩短；其次，基金倾向于追求保值增值导致偏向成熟项目，偏离“投早投小”的目标。引导基金的政策初衷是支持处于早期阶段、创新含量高但风险也高的科技企业。在债务压力下，地方政府对财政资金保值增值的要求更为严格，国有资产考核体系使得管理人“不敢投”初创期企业，宁愿投向确定性高但支持必要性不高的成熟期项目，甚至购买理财产品。

考核体系的变形使得地方政府倾向于缩短考核周期，要求基金尽快呈现正向回报。田轩指出，\erjfootnote{北京大学博雅特聘教授田轩在2026博鳌亚洲论坛年会上的相关发言。}国资考核机制与耐心资本的市场化表现存在矛盾——创投容错机制不健全，收益负面波动时缺乏包容甚至问责处罚。当考核周期被压缩，基金管理人不再有“耐心等待”创新产出的空间，投资决策从“选择最有长期价值的创新项目”转向“选择能在考核期内呈现正回报的安全项目”。审计发现，山西8.12亿元基金违规投向房地产开发等项目，\erjfootnote{山西省审计厅《关于2023年度省级预算执行和其他财政收支的审计工作报告》。}云南某公司借用政府投资基金1.54亿元用于偿还债务。\erjfootnote{云南省人大常委会2023年度审计查出问题整改情况报告（2024年11月发布）。}当债务压力较低时，地方政府能够容忍引导基金的合理亏损，投资决策向“投早投小”倾斜；当债务压力较高时，投资决策向“保本保息”倾斜，创新最需要支持的早期科技企业反而得不到资金。Zhao和Zhang（2025）的研究也证实了这一点——引导基金在高新技术产业和创投成熟的地区对创新边界的拓展效应更强，而这些地区恰恰是债务压力相对较低的地区。最后，值得考虑的是，在债务压力下，行政干预加剧会造成投资决策变形，进而侵蚀耐心资本的属性。面临债务压力的地方政府，更可能将引导基金视为缓解财政紧张或实现短期政绩的工具，加强行政干预。审计显示，宁夏4400万元基金因要求的投资收益率过高导致3520万元闲置，\erjfootnote{宁夏2023年度重大政策落实及重大项目审计情况（2024年8月发布）。}辽宁1只基金到位5.02亿元仅投资4830万元（占比9.66\%）。\erjfootnote{辽宁省审计厅相关审计报告（2024年发布）。}Wang等（2025）发现政府干预较少、财政压力较低的地区，债务对创新质量的负面影响更小，耐心资本属性发挥的作用会更强。

由此提出假说H3：债务压力通过削弱引导基金的耐心资本属性，进一步抑制其创新促进作用。

政府产业引导基金能以少量财政资金撬动数倍社会资本，形成规模化创新资本供给（Xiao et al., 2024；蔡庆丰等，2024）。然而，地方政府债务压力从多个维度削弱了这一撬动效率，抑制了企业创新。

首先，债务压力会挤占财政空间，导致引导基金出资额减少（崔竹，2025）。当地方政府债务压力较低时，财政资源相对充裕，政府能够按计划足额实缴承诺出资，引导基金得以顺利运作。当债务压力较高时，偿债支出和“三保”等刚性支出挤占财政空间，政府可用于引导基金的财政资金被严重压缩，社会资本的募集减少；其次是债务压力会诱致“名股实债”等异化行为，即引导基金从风险投资工具异化为债权工具。为弥补社会资本信心不足，债务压力大的地区更倾向于通过承诺回购本金、设定最低投资收益或承担投资损失等方式吸引社会资本（崔军等，2025）。社会资本的“撬动”功能被“替代”功能所取代。罗志恒（2021）测算，引导基金通过杠杆承诺、劣后级兜底等机制使地方政府隐性债务增加约1.96至4.54万亿元，形成“债务越高—越需要保底募资—隐性债务越膨胀—创新功能越弱”的恶性循环；此外，社会资本对高债务地区信心下降，募资难度加大，直接降低了引导基金对社会资本的撬动效率（张圆圆、夏球，2024）。当债务压力较低时，地方财政实力强、信用好，社会资本参与意愿高；当债务压力较高时，社会资本对地方财政信心下降，社会资本在高债务地区的投资意愿显著下降，企业创新获得的总资本量大幅减少，撬动效率被严重削弱。撬动效率的降低直接限制了可用于支持企业创新的资金规模，企业创新的融资约束难以得到有效缓解（Hao et al., 2025）。

由此提出假说H4:债务压力通过降低社会资本撬动效率，进一步抑制基金创新促进作用。

初创期和早中期的科技型企业面临三重典型的市场失灵：信息不对称、正外部性、缺乏用于融资的可抵押资产，三重作用叠加使得其难以仅靠市场机制作用与企业自身管理缓解融资约束。相较于依赖稳定现金流和抵押资产的债权融资，股权融资更能分担科技创新活动中的高风险与长周期特征，也更契合科技型企业持续研发投入的资金需求。政府产业引导基金根据企业所处生命周期的不同阶段，实施阶段性投资与联合投资，匹配企业差异化的资金需求，直接为企业创新活动提供长期资本支持。同时，政府产业引导基金的发展有利于促进地区金融资源的均衡分布，资本资源聚集的网络效应有利于企业信息在资本之间的共享，降低社会资本寻求投资对象过程中的信息成本（李健旋和赵林度，2018）。因此，政府产业引导基金不仅能够通过股权投资直接缓解科技型企业的资金缺口，还能够通过政府背书、风险分担与联合投资向社会资本释放质量认证信号，进而提升企业获得外部股权融资的可能性。

然而地方债务压力的攀升从源头上削弱产业引导基金缓解企业融资约束的能力。地方政府债务过高源于政府投资性项目偏离实际公共需求，从而造成政府支出低效率配置，不利于当地经济增长（陈思霞和陈志勇，2015)。在此背景下，债务率较高地区的地方财政逐渐进入“紧平衡”状态。在经济下行压力下，地方政府面临融资渠道受限与债务风险上升等问题（崔竹，2025），财政空间进一步收窄，从而压缩了政府对产业引导基金持续出资的能力。我国各级政府投资基金的累计认缴率约为64.9\%\footnote[1]{\hspace*{2em}数据来源于清科研究中心. 2025年中国政府引导基金发展研究报告.}，大量已认缴资金无法实际到位，直接影响引导基金的正常运作。更重要的是，地方债务压力会削弱政府产业引导基金作为股权融资工具的风险承担能力与资本撬动能力。一方面，财政紧平衡会使地方政府更倾向于将有限资金投向短期稳增长和债务化解任务，降低对初创期、早中期科技型企业这类高风险、长周期项目的持续支持力度；另一方面，财政出资不到位和政策连续性下降会削弱政府资金的信用背书功能，降低社会资本对引导基金的跟投意愿，使基金原本通过联合投资放大资本供给、缓解信息不对称的机制难以充分发挥。除了直接影响资金来源外，地方债务压力还会通过压缩整体融资环境进一步削弱政府投资基金的作用。研究指出，地方政府债务扩张会通过直接占压银行信贷资源，对企业融资形成“挤出效应”，这种挤出效应在“中小企业和非国有企业”中表现得比大型企业和国有企业更为显著（华夏等，2020）。我国金融市场银行信贷的配置效率较低下，在缺乏融资的抵押资产初创企业中，银行更加青睐有政府信用作为隐性担保的国有企业和融资平台（周天宇，2025）。当引导基金自身出资能力受限、社会资本撬动效率下降，而外部融资环境又因债务扩张而趋于收紧时，初创期和早中期科技型企业面临更强的股权融资约束，其融资约束被进一步恶化，从而严重制约研发投入与持续创新能力的提升。

由此提出假说H5：债务压力通过弱化基金缓解融资约束的能力，进一步抑制城市创新。
% ===== END INPUT: 二_文献综述与理论假说.tex =====



% ===== BEGIN INPUT: 三_模型设定与变量说明.tex =====
% 本文件由 main.tex 输入，作为“模型设定与变量说明”章节，不单独编译。

\section{模型设定与变量说明}

\subsection{数据来源与样本说明}

本文以地级市—年份为研究单元，样本期为2015—2024年。选择地级市层面展开分析，主要是因为政府产业引导基金通常以注册地、出资地和服务地方产业为重要运行边界，地方债务压力和财政约束也主要在地方政府层面体现。与企业层面数据相比，地级市面板更适合识别地方财政债务约束下政府产业引导基金与区域创新产出之间的关系。

本文使用的创新产出数据来自 CNRDS 地方专利数据，并整理为地级市—年份层面的发明申请量、实用新型申请量和专利申请总量。政府产业引导基金基础信息来自清科基金目录和投中基金目录，本文将两类目录合并去重后，依据基金注册地手工匹配到地级市层面，进而构造各城市基金当年设立数量、累计设立数量和累计设立规模。基金出资结构主要来自投中基金目录中的 LP 和出资人信息，用以识别政府资本、国资平台、政府产业引导基金、FOF 以及其他社会资本出资。基金投资事件信息主要来自清科投资事件数据，本文据此识别政府引导基金的投资时间、投资阶段、投资次数和投资金额，并按基金注册地汇总到地级市—年份层面。地方债务数据依据各地方统计局和统计年鉴披露信息手工整理，控制变量主要来自 CEIC、CNKI 等城市统计资料。

样本处理遵循现有回归面板的统一口径。回归前按城市和年份进行去重，模型中加入城市固定效应和年份固定效应，并按城市维度聚类稳健标准误。由于部分债务变量、基金投资事件变量和机制变量存在缺失，不同回归表中的样本量会随变量口径变化，本文在实证结果表中分别报告各规格观测值数量。

\subsection{变量定义}

本文被解释变量为城市创新产出。正文主结果使用专利申请量而非授权量，原因在于申请量能够及时地反映城市创新活动变化。具体而言，本文分别使用发明申请量、实用新型申请量和专利申请总量衡量创新产出，其中发明申请量更接近高质量创新口径，专利申请总量则作为综合创新产出口径。涉及对数形式时，统一使用 $\ln(1+\text{专利量})$ 处理。

核心解释变量为政府产业引导基金规模。基准回归同时考察基金当年设立数量、累计设立数量和累计设立规模。相较于基金数量口径，累计设立规模更直接反映某地截至当年通过引导基金形成的政策性资本供给总量，也更贴近本文关注的政府资本支持强度和资源配置能力，因此本文将其作为正文主线变量，并将基金当年设立数量和累计设立数量作为替代基金口径。

调节变量为地方债务压力。本文采用债券余额合计与公共财政收入之比衡量地方债务压力，该值越高，表示地方偿债压力和财政约束越强。相较于财政收支压力等流量指标，该指标更直接刻画地方政府存量债务负担及其对财政空间、风险承担能力的约束。为缓解同期反向因果和共同冲击的影响，正文主规格优先采用滞后一期债务压力；当期债务压力作为补充口径。财政收支压力更多反映年度预算平衡紧张程度，因此在本文中仅作补充讨论，不作为正文主调节变量。

机制变量围绕三条理论路径设置。第一，耐心资本属性用早期投资事件占比衡量。耐心资本的核心在于能够承受较高不确定性和较长回报周期，并将资金配置到短期现金流不确定但创新潜力较高的项目。张杰和任元明（2025）指出，耐心资本能够为企业自主创新和产业链培育提供跨周期资金支持；张壹帆和陆岷峰（2025）也将政府引导基金“耐心度”与投早中期项目数量、长期绩效评价和容错机制联系起来。政府风险资本文献同样把种子期、初创期或早期企业投资作为识别公共资本是否补足早期创新融资缺口的重要口径（Cumming，2007；Cumming and Johan，2007）。在国内情境下，施国平等（2016）直接使用创投机构投资早期创业企业轮次占总投资轮次的比例，检验政府引导基金是否引导创投机构投向早期企业。因此，本文将投资阶段为种子期或初创期的事件记为早期投资，早期投资事件占比等于早期投资事件数除以该市引导基金当年投资总次数。相较于早期投资金额占比，事件占比更直接反映基金是否将投资机会配置到早期项目，也较少受少数大额投资事件影响。第二，社会资本撬动效率用社会资本出资与政府资本出资之比及其变化衡量，其中社会资本出资剔除国资平台、政府产业引导基金和 FOF 后识别。第三，融资约束缓解效应借鉴王平、徐肇仪（2023）将金融结构作为融资约束代理变量的思路，使用存款/GDP口径的金融发展水平作为融资约束的反向代理。该指标越高，表示地区金融资源供给和吸纳能力越强，地区融资约束相对越弱。需要说明的是，该变量衡量的是城市层面的融资约束缓解程度，不等同于企业层面的 SA、KZ 或 WW 融资约束指数。

\begin{table}[htbp]
  \centering
  \caption*{主要变量定义与构造}
  {\ERJTableFont
  \begin{tabularx}{0.96\linewidth}{l l X}
    \toprule
    变量类型 & 变量名称 & 构造说明 \\
    \midrule
    创新产出 & 发明申请量 & CNRDS 地方专利数据中地级市—年份发明专利申请数量，用于衡量高质量创新产出。 \\
    创新产出 & 实用新型申请量 & CNRDS 地方专利数据中地级市—年份实用新型专利申请数量，用作补充创新口径。 \\
    创新产出 & 专利申请总量 & 发明、实用新型和外观设计等专利申请数量汇总，用于衡量综合创新产出。 \\
    基金变量 & 基金当年设立数量 & 清科和投中基金目录合并去重后，按注册地和设立年份汇总到城市层面。 \\
    基金变量 & 基金累计设立数量 & 截至当年各地级市累计设立的政府产业引导基金数量。 \\
    基金变量 & 基金累计设立规模 & 截至当年各地级市累计设立的政府产业引导基金规模，正文主解释变量。 \\
    债务压力 & 债务压力 & 债券余额合计除以公共财政收入，数值越高表示地方债务约束越强。 \\
    债务压力 & 滞后一期债务压力 & 债务压力的一期滞后项，正文主调节变量。 \\
    耐心资本 & 早期投资事件占比 & 种子期和初创期投资事件数占该市引导基金当年投资事件总数的比例。 \\
    社会资本 & 社会资本撬动效率 & 非政府社会资本出资与政府资本出资之比，用于衡量政府资本放大效应。 \\
    融资约束 & 存款/GDP & 年末金融机构存款余额与地区生产总值之比，用作地区融资约束的反向代理。 \\
    \bottomrule
  \end{tabularx}

  \par\smallskip
  {\ERJTableNoteFont 注：涉及对数变量时，规模变量统一采用 $\ln(x+1)$ 处理；融资约束机制使用存款/GDP口径金融发展水平作为反向代理。}
  }
\end{table}

表1报告主要变量的描述性统计。样本中城市专利申请总量均值为13109.508件，发明申请量均值为4683.297件。基金累计设立规模均值为4.439亿元，但标准差达到16.662亿元，说明政府产业引导基金在城市之间分布差异较大。债务压力均值为233.176，滞后一期债务压力均值为199.112，也显示地方债务约束存在明显地区差异。早期投资事件占比均值为0.340，表明政府引导基金投资事件中约三分之一投向种子期或初创期项目。

\begin{table}[htbp]
  \centering
  \caption{主要变量描述性统计}
  {\ERJTableFont
  \begin{tabular}{@{}L{4.2cm} r r r r r@{}}
    \toprule
    变量 & N & 均值 & 标准差 & 最小值 & 最大值 \\
    \midrule
    发明申请量 & 2720 & 4683.297 & 13016.643 & 6.000 & 195388.000 \\
    实用新型申请量 & 2720 & 6420.434 & 13622.132 & 21.000 & 146592.000 \\
    专利申请总量 & 2720 & 13109.508 & 29764.731 & 34.000 & 323014.000 \\
    基金当年设立数量 & 2770 & 2.803 & 7.215 & 0.000 & 99.000 \\
    基金累计设立数量 & 2770 & 14.266 & 39.521 & 0.000 & 458.000 \\
    基金累计设立规模（亿元） & 2629 & 4.439 & 16.662 & 0.000 & 279.266 \\
    债务压力 & 1904 & 233.176 & 484.577 & 0.746 & 4654.410 \\
    滞后一期债务压力 & 1898 & 199.112 & 422.855 & 0.806 & 4129.294 \\
    早期投资事件占比 & 957 & 0.340 & 0.338 & 0.000 & 1.000 \\
    社会资本撬动效率 & 986 & 0.789 & 2.015 & 0.000 & 34.050 \\
    金融发展水平 & 2500 & 1.638 & 0.772 & 0.000 & 20.100 \\
    经济发展水平 & 2098 & 5.453 & 0.864 & 3.081 & 8.590 \\
    财政科技支出 & 1971 & 6.343 & 1.418 & 2.757 & 11.060 \\
    人口规模 & 2450 & 5.980 & 0.748 & 3.252 & 11.975 \\
    第二产业规模 & 2098 & 4.554 & 0.886 & 2.272 & 7.238 \\
    外资利用水平 & 1681 & 10.184 & 2.082 & 1.386 & 14.776 \\
    \bottomrule
  \end{tabular}

  \par\smallskip
  {\ERJTableNoteFont 注：基金累计设立规模按亿元报告。经济发展水平、财政科技支出、人口规模、第二产业规模和外资利用水平为 $\ln(x+1)$ 口径。社会资本撬动效率报告有限值观测，剔除因政府出资为0形成的非有限值。}
  }
\end{table}

控制变量沿用现有回归设定，主要包括经济发展水平、财政科技支出、人口规模、产业结构和外资利用水平。经济发展水平用地区生产总值对数衡量，财政科技支出用于控制地方财政对科技创新的直接支持强度，常住人口对数用于控制人口规模和市场规模，第二产业增加值对数用于控制城市产业基础，实际利用外资额对数用于控制对外开放程度。上述控制变量均按 $\ln(x+1)$ 处理，以缓解规模变量右偏分布并保留取值为零的观测。

\subsection{模型设定}

为检验政府产业引导基金是否促进城市创新，本文首先估计如下双向固定效应模型：

\begin{equation}
Innovation_{it}=\alpha+\beta Fund_{it}+\gamma Controls_{it}+\mu_i+\lambda_t+\varepsilon_{it}.
\end{equation}
\ERJWhere{$Innovation_{it}$ 表示城市 $i$ 在年份 $t$ 的创新产出，$Fund_{it}$ 表示政府产业引导基金变量，$Controls_{it}$ 为城市层面控制变量，$\mu_i$ 和 $\lambda_t$ 分别表示城市固定效应和年份固定效应。若 $\beta$ 显著为正，则说明政府产业引导基金扩张与城市创新产出提升相一致。}

在此基础上，本文引入债务压力与基金规模的交互项，考察地方债务压力是否削弱基金的创新促进效应：

\begin{equation}
Innovation_{it}=\alpha+\beta_1 Fund_{it}+\beta_2 Debt_{i,t-1}
+\beta_3(Fund_{it}\times Debt_{i,t-1})+\gamma Controls_{it}+\mu_i+\lambda_t+\varepsilon_{it}.
\end{equation}
\ERJWhere{$Debt_{i,t-1}$ 表示滞后一期债务压力，$Fund_{it}\times Debt_{i,t-1}$ 为核心交互项。本文关注 $\beta_3$ 的符号和显著性。若 $\beta_3<0$，说明债务压力越高，同等规模政府产业引导基金对创新产出的边际促进作用越弱。}

机制检验围绕耐心资本属性、社会资本撬动效率和融资约束缓解效应展开。耐心资本和融资约束机制主要采用“债务压力削弱基金影响机制变量，机制变量进一步影响创新产出”的路径检验：

\begin{equation}
M_{it}=\alpha+\theta_1(Fund_{it}\times Debt_{i,t-1})+\gamma Controls_{it}+\mu_i+\lambda_t+\nu_{it},
\end{equation}

\begin{equation}
Innovation_{i,t+k}=\alpha+\rho_1(Fund_{it}\times Debt_{i,t-1})+\rho_2 M_{it}
+\gamma Controls_{it}+\mu_i+\lambda_t+\eta_{it}.
\end{equation}
\ERJWhere{$M_{it}$ 表示机制变量，$k$ 表示创新产出的滞后实现期。耐心资本机制中 $M_{it}$ 为早期投资事件占比，理论预期为 $\theta_1<0$ 且 $\rho_2>0$；融资约束机制中 $M_{it}$ 为存款/GDP衡量的金融发展水平，理论预期同样为 $\theta_1<0$ 且 $\rho_2>0$。}

社会资本撬动效率机制不写作普通中介，而写作调节式机制或承接机制。相应模型在结果方程中进一步加入基金规模与社会资本撬动效率的交互项：

\begin{equation}
Innovation_{it}=\alpha+\delta_1(Fund_{it}\times Debt_{i,t-1})
+\delta_2(Fund_{it}\times SocCap_{it})+\gamma Controls_{it}+\mu_i+\lambda_t+\xi_{it}.
\end{equation}
\ERJWhere{$SocCap_{it}$ 表示社会资本撬动效率或其变化。若债务压力降低撬动效率，同时 $\delta_2>0$，且加入机制项后 $\delta_1$ 的绝对值下降，则说明债务压力对基金创新效应的抑制部分通过社会资本撬动效率这一承接渠道体现。}

以上模型均采用城市固定效应和年份固定效应，并按城市维度聚类稳健标准误。本文的识别重点不是宣称完全解决所有内生性问题，而是在固定城市不随时间变化的特征、年份共同冲击以及主要城市层面控制变量后，比较同一城市在不同年份基金规模、债务压力和创新产出的共同变化关系。后续滞后债务压力和滞后创新响应的检验，用于进一步缓解反向因果和同期冲击的担忧。
% ===== END INPUT: 三_模型设定与变量说明.tex =====



% ===== BEGIN INPUT: 四_实证结果分析.tex =====
% 本文件由 main.tex 输入，作为“实证结果分析”章节，不单独编译。

\section{实证结果分析}

\subsection{基准回归：政府产业引导基金与城市创新}

政府产业引导基金扩张与城市创新产出显著正相关。表2显示，在双向固定效应设定下，基金累计设立规模对发明申请量、实用新型申请量和专利申请总量的系数分别为0.030、0.041和0.080，均在常规水平上显著。替换为基金当年设立数量或累计设立数量后，核心结论仍保持一致。

\begin{table}[htbp]
  \centering
  \caption{政府产业引导基金与城市创新：基准回归}
  {\ERJTableFont
  \begin{tabular}{@{}L{2.8cm} L{2.8cm} C{3.0cm} C{1.1cm} C{1.1cm}@{}}
    \toprule
    $Innovation_{it}$ & $Fund_{it}$口径 & 系数（标准误） & N & $R^2$ \\
    \midrule
    发明申请量 & 累计设立规模 & \coefse{0.030***}{0.009} & 1083 & 0.462 \\
    实用新型申请量 & 累计设立规模 & \coefse{0.041***}{0.014} & 1083 & 0.516 \\
    专利申请总量 & 累计设立规模 & \coefse{0.080**}{0.033} & 1083 & 0.495 \\
    发明申请量 & 当年设立数量 & \coefse{297.19***}{86.54} & 1149 & 0.203 \\
    专利申请总量 & 当年设立数量 & \coefse{849.60***}{267.97} & 1149 & 0.296 \\
    发明申请量 & 累计设立数量 & \coefse{106.80***}{30.85} & 1149 & 0.449 \\
    \midrule
    $Controls_{it}$ & \multicolumn{4}{c}{是} \\
    $\mu_i$ & \multicolumn{4}{c}{是} \\
    $\lambda_t$ & \multicolumn{4}{c}{是} \\
    \bottomrule
  \end{tabular}

  \par\smallskip
  {\ERJTableNoteFont 注：括号内为按城市聚类的稳健标准误。*、**、*** 分别表示在10\%、5\%、1\% 水平上显著。}
  }
\end{table}

表2支持假说 H1。累计设立规模对发明申请量显著为正，说明基金扩张不仅对应一般专利数量增长，也对应更接近高质量创新的专利申请增加。基金数量口径同样为正，表明结论不完全依赖基金规模变量。

\subsection{债务压力的调节效应}

地方债务压力显著削弱基金的创新促进效应。表3承接前文调节模型，报告基金累计设立规模与滞后一期债务压力的交互项结果。交互项在专利申请总量、发明申请量和实用新型申请量三个口径下均显著为负，系数分别为 $-6.09\times 10^{-5}$、$-1.62\times 10^{-5}$ 和 $-2.92\times 10^{-5}$。

\begin{table}[htbp]
  \centering
  \caption{债务压力对基金创新效应的调节作用}
  {\ERJTableFont
  \begin{tabular}{@{}L{3.1cm} C{4.4cm} C{1.1cm} C{1.1cm}@{}}
    \toprule
    $Innovation_{it}$ & $Fund_{it}\times Debt_{i,t-1}$ 系数（标准误） & N & $R^2$ \\
    \midrule
    专利申请总量 & \coefse{$-6.09\times 10^{-5}$***}{$1.66\times 10^{-5}$} & 914 & 0.693 \\
    发明申请量 & \coefse{$-1.62\times 10^{-5}$***}{$4.43\times 10^{-6}$} & 914 & 0.570 \\
    实用新型申请量 & \coefse{$-2.92\times 10^{-5}$***}{$4.88\times 10^{-6}$} & 914 & 0.691 \\
    \midrule
    $Controls_{it}$ & \multicolumn{3}{c}{是} \\
    $\mu_i$ & \multicolumn{3}{c}{是} \\
    $\lambda_t$ & \multicolumn{3}{c}{是} \\
    \bottomrule
  \end{tabular}

  \par\smallskip
  {\ERJTableNoteFont 注：交互项为 $Fund_{it}\times Debt_{i,t-1}$，其中 $Fund_{it}$ 采用基金累计设立规模，$Debt_{i,t-1}$ 为滞后一期债务压力。括号内为按城市聚类的稳健标准误。财政压力结果不作为正文主调节结论。}
  }
\end{table}

表3说明，同等规模的政府产业引导基金在高债务压力城市中更难转化为创新产出。一个合理解释是，债务约束压缩地方政府财政空间和风险承担能力，使基金更容易受到短期绩效和偿债目标约束。该结果支持 H2：地方债务压力会削弱政府产业引导基金的创新促进效应。

\subsection{稳健性检验：替换变量口径}

表4检验结论是否依赖单一变量口径。替换创新产出、替换基金变量、将发明申请量改为对数形式后，基金变量仍为正；将滞后一期债务压力替换为当期债务压力后，核心交互项仍在三类创新产出下显著为负。

\begin{table}[htbp]
  \centering
  \caption{稳健性检验：替换变量口径}
  {\ERJTableFont
  \begin{tabular}{@{}L{2.0cm} L{2.3cm} L{3.4cm} C{2.9cm} C{0.95cm} C{0.95cm}@{}}
    \toprule
    检验内容 & $Innovation_{it}$ & 核心项 & 系数（标准误） & N & $R^2$ \\
    \midrule
    替换创新口径 & 实用新型申请量 & $Fund_{it}$（当年设立数量） & \coefse{458.08***}{125.55} & 1149 & 0.367 \\
    替换基金口径 & 专利申请总量 & $Fund_{it}$（累计设立数量） & \coefse{296.77**}{116.95} & 1149 & 0.513 \\
    替换因变量形式 & 发明申请量对数 & $Fund_{it}$（累计设立规模） & \coefse{$4.02\times 10^{-7}$**}{$1.72\times 10^{-7}$} & 1083 & 0.339 \\
    替换债务口径 & 专利申请总量 & $Fund_{it}\times Debt_{it}$ & \coefse{$-5.56\times 10^{-5}$***}{$1.41\times 10^{-5}$} & 917 & 0.707 \\
    替换债务口径 & 发明申请量 & $Fund_{it}\times Debt_{it}$ & \coefse{$-1.45\times 10^{-5}$***}{$3.45\times 10^{-6}$} & 917 & 0.586 \\
    替换债务口径 & 实用新型申请量 & $Fund_{it}\times Debt_{it}$ & \coefse{$-2.74\times 10^{-5}$***}{$4.33\times 10^{-6}$} & 917 & 0.696 \\
    \bottomrule
  \end{tabular}

  \par\smallskip
  {\ERJTableNoteFont 注：所有规格均加入 $Controls_{it}$、$\mu_i$ 和 $\lambda_t$，并按城市聚类稳健标准误。}
  }
\end{table}

表4表明，基准结论和债务压力调节结论不是由某一个变量口径单独驱动。对数因变量结果弱于水平值口径，但发明申请量对数仍提供了方向一致的支持。债务压力采用当期口径时，负向调节效应也没有消失。

\subsection{内生性处理：滞后变量检验}

表5比较当期债务压力和滞后一期债务压力两类口径。两类设定下，基金累计设立规模与债务压力的交互项都显著为负，且滞后一期债务压力的结果与当期口径接近。

\begin{table}[htbp]
  \centering
  \caption{内生性处理：滞后债务压力检验}
  {\ERJTableFont
  \begin{tabular}{@{}L{2.5cm} L{2.8cm} C{4.0cm} C{1.0cm} C{1.0cm}@{}}
    \toprule
    $Innovation_{it}$ & $Debt$口径 & $Fund_{it}\times Debt$ 系数 & N & $R^2$ \\
    \midrule
    专利申请总量 & $Debt_{it}$ & \coefse{$-5.56\times 10^{-5}$***}{$1.41\times 10^{-5}$} & 917 & 0.707 \\
    专利申请总量 & $Debt_{i,t-1}$ & \coefse{$-6.09\times 10^{-5}$***}{$1.66\times 10^{-5}$} & 914 & 0.693 \\
    发明申请量 & $Debt_{it}$ & \coefse{$-1.45\times 10^{-5}$***}{$3.45\times 10^{-6}$} & 917 & 0.586 \\
    发明申请量 & $Debt_{i,t-1}$ & \coefse{$-1.62\times 10^{-5}$***}{$4.43\times 10^{-6}$} & 914 & 0.570 \\
    实用新型申请量 & $Debt_{it}$ & \coefse{$-2.74\times 10^{-5}$***}{$4.33\times 10^{-6}$} & 917 & 0.696 \\
    实用新型申请量 & $Debt_{i,t-1}$ & \coefse{$-2.92\times 10^{-5}$***}{$4.88\times 10^{-6}$} & 914 & 0.691 \\
    \bottomrule
  \end{tabular}

  \par\smallskip
  {\ERJTableNoteFont 注：交互项为 $Fund_{it}\times Debt$，其中 $Fund_{it}$ 采用基金累计设立规模。所有规格均加入 $Controls_{it}$、$\mu_i$ 和 $\lambda_t$，并按城市聚类稳健标准误。}
  }
\end{table}

表5不能完全排除所有内生性问题，但它降低了结论由同期创新冲击驱动的可能性。后续机制检验进一步使用未来创新产出，以允许基金投资结构对专利申请产生滞后反应。

\subsection{机制检验一：耐心资本属性}

第一条机制关注基金的耐心资本属性。按照第三章的变量定义，早期投资事件占比用于刻画基金是否真正配置到风险更高、周期更长的早期项目。若债务压力压缩基金的长期配置能力，则交互项应降低早期投资事件占比，而早期投资事件占比应提高后续创新产出。为避免比例分母过小导致机械波动，表6将样本限定为基金投资事件数不少于2的城市—年份观测。

\begin{table}[htbp]
  \centering
  \caption{耐心资本属性机制：早期投资事件占比}
  {\ERJTableFont
  \resizebox{\linewidth}{!}{%
  \begin{tabular}{@{}l l l c c c c c c@{}}
    \toprule
    $M_{it}$口径 & $Innovation_{i,t+k}$ & $Debt$口径 & $\beta_3$（基准） & $\theta_1$ & $\rho_2$ & $\rho_1$（加 $M_{it}$ 后） & N & $R^2$ \\
    \midrule
    原始比例 & 专利申请总量（$k=2$） & $Debt_{i,t-1}$ & \coefse{-0.085**}{0.041} & \coefse{-0.143**}{0.065} & \coefse{0.092**}{0.045} & \coefse{-0.071*}{0.042} & 266 & 0.680 \\
    比例1/99缩尾 & 专利申请总量（$k=2$） & $Debt_{i,t-1}$ & \coefse{-0.085**}{0.041} & \coefse{-0.143**}{0.065} & \coefse{0.092**}{0.045} & \coefse{-0.071*}{0.042} & 266 & 0.680 \\
    原始比例 & 专利申请总量（$k=2$） & $Debt_{i,t-1}$（1/99缩尾） & \coefse{-0.075**}{0.036} & \coefse{-0.128**}{0.057} & \coefse{0.091**}{0.045} & \coefse{-0.064*}{0.037} & 266 & 0.680 \\
    比例1/99缩尾 & 专利申请总量（$k=2$） & $Debt_{i,t-1}$（1/99缩尾） & \coefse{-0.075**}{0.036} & \coefse{-0.128**}{0.057} & \coefse{0.091**}{0.045} & \coefse{-0.064*}{0.037} & 266 & 0.680 \\
    \bottomrule
  \end{tabular}
  }

  \par\smallskip
  {\ERJTableNoteFont 注：$\beta_3$ 表示未加入 $M_{it}$ 时 $Fund_{it}\times Debt_{i,t-1}$ 对未来创新的影响；$\theta_1$ 表示 $Fund_{it}\times Debt_{i,t-1}$ 对 $M_{it}$ 的影响；$\rho_2$ 表示 $M_{it}$ 对未来创新的影响；$\rho_1$ 为加入 $M_{it}$ 后结果方程中继续保留的 $Fund_{it}\times Debt_{i,t-1}$ 系数。括号内为按城市聚类的稳健标准误。所有规格均加入 $Controls_{it}$、$\mu_i$ 和 $\lambda_t$。}
  }
\end{table}

表6显示，债务压力显著降低早期投资事件占比，而早期投资事件占比显著提高未来两年专利申请总量。加入早期投资事件占比后，基金规模与债务压力交互项仍为负，但绝对值和显著性均有所下降。该结果支持 H3：债务压力可能通过削弱引导基金的耐心资本属性，进一步抑制其创新促进作用。该机制证据的边界在于，结果来自分母稳定化样本，早期投资金额占比口径并不稳定。

\subsection{机制检验二：社会资本撬动效率}

第二条机制关注社会资本撬动效率。按照前文的调节式承接机制设定，表7重点检验两点：债务压力是否削弱撬动效率，以及较高撬动效率是否强化基金规模对创新的边际促进作用。
\begin{table}[htbp]
  \centering
  \caption{社会资本撬动效率机制：低金融发展地区}
  {\ERJTableFont
  \resizebox{\linewidth}{!}{%
  \begin{tabular}{@{}l l c c c c c c@{}}
    \toprule
    $Innovation_{i,t+k}$ & $Debt$口径 & 基准 $Fund_{it}\times Debt$ & $SocCap_{it}$方程 & $Fund_{it}\times SocCap_{it}$ & 加机制后 $Fund_{it}\times Debt$ & N & $R^2$ \\
    \midrule
    下一期实用新型申请量 & $Debt_{i,t-1}$ & \coefse{$-6.82\times 10^{-4}$*}{$3.72\times 10^{-4}$} & \coefse{$-7.89\times 10^{-8}$*}{$4.39\times 10^{-8}$} & \coefse{0.065***}{0.013} & \coefse{$-6.03\times 10^{-4}$*}{$3.43\times 10^{-4}$} & 243 & 0.569 \\
    下一期专利申请总量 & $Debt_{i,t-1}$ & \coefse{$-1.37\times 10^{-3}$**}{$5.45\times 10^{-4}$} & \coefse{$-7.89\times 10^{-8}$*}{$4.39\times 10^{-8}$} & \coefse{0.081***}{0.019} & \coefse{$-1.26\times 10^{-3}$**}{$5.10\times 10^{-4}$} & 243 & 0.472 \\
    \midrule
    实用新型申请量 & $Debt_{i,t-1}$ & \coefse{-0.014***}{0.005} & \coefse{-0.094*}{0.056} & \coefse{0.113*}{0.058} & \coefse{-0.012***}{0.004} & 121 & 0.749 \\
    下一期专利申请总量 & $Debt_{i,t-1}$ & \coefse{-0.028***}{0.010} & \coefse{-0.094*}{0.056} & \coefse{0.182*}{0.100} & \coefse{-0.025***}{0.008} & 121 & 0.606 \\
    \bottomrule
  \end{tabular}
  }

  \par\smallskip
  {\ERJTableNoteFont 注：表中结果均来自低金融发展样本。前两行的 $SocCap_{it}$ 方程对应 $Fund_{it}\times Debt_{i,t-1}\rightarrow SocCap_{it}$；后两行为加入经济发展水平、财政科技支出、人口规模、产业结构和外资利用水平后的10\% 显著结果，$SocCap_{it}$ 方程对应 $Debt_{i,t-1}\rightarrow SocCap_{it}$。机制变量均为剔除极小政府出资后缩尾的社会资本撬动效率。$R^2$ 为结果方程拟合优度。加入机制项后，原 $Fund_{it}\times Debt$ 项绝对值下降。}
  }
\end{table}

表7显示，在金融发展水平较低的城市中，债务压力会削弱社会资本撬动效率；撬动效率越高，基金规模对创新产出的边际促进作用越强。加入机制项后，原基金规模与债务压力交互项的绝对值下降。该结果表示上述机制在低金融发展地区成立，从而使H4在低金融发展水平地区受到支持。

\subsection{机制检验三：融资约束缓解效应}

第三条机制检验融资约束缓解效应。本文借鉴王平、徐肇仪（2023）将金融结构作为融资约束代理变量的思路，并与前文金融发展口径保持一致，使用存款/GDP衡量地区金融发展水平，作为融资约束的反向代理。该指标越高，表示地区金融资源供给和吸纳能力越强，融资约束相对越弱。为保证完整机制链条成立，表8进一步限定基准调节项、机制方程交互项和结果方程机制变量均满足理论方向，其中受控模型下结果方程机制变量至少达到10\%显著水平。

\begin{table}[htbp]
  \centering
  \caption{融资约束缓解机制：严格筛选结果}
  {\ERJTableFont
  \resizebox{\linewidth}{!}{%
  \begin{tabular}{@{}l l l c c c c c c@{}}
    \toprule
    $M_{it}$口径 & $Innovation_{i,t+k}$ & $Debt$口径 & $\beta_3$（基准） & $\theta_1$ & $\rho_2$ & $\rho_1$（加 $M_{it}$ 后） & N & $R^2$ \\
    \midrule
    存款/GDP增速 & 专利申请总量 & $Debt_{it}$ & \coefse{-6.244***}{1.892} & \coefse{$-5.77\times 10^{-6}$**}{$2.24\times 10^{-6}$} & \coefse{11875.512*}{6385.986} & \coefse{-6.176***}{1.871} & 784 & 0.566 \\
    存款/GDP增速 & 实用新型申请量 & $Debt_{it}$ & \coefse{-3.517***}{0.866} & \coefse{$-5.77\times 10^{-6}$**}{$2.24\times 10^{-6}$} & \coefse{6727.242*}{3792.550} & \coefse{-3.478***}{0.852} & 784 & 0.597 \\
    存款/GDP（asinh变换） & 专利申请总量对数 & $\ln(1+Debt_{i,t-1})$ & \coefse{-0.0069**}{0.0032} & \coefse{-0.0018**}{0.0009} & \coefse{0.341*}{0.197} & \coefse{-0.0063*}{0.0032} & 944 & 0.706 \\
    存款/GDP（对数变换） & 专利申请总量对数 & $\ln(1+Debt_{i,t-1})$ & \coefse{-0.0069**}{0.0032} & \coefse{-0.0013**}{0.0006} & \coefse{0.455*}{0.267} & \coefse{-0.0063*}{0.0032} & 944 & 0.706 \\
    \bottomrule
  \end{tabular}
  }

  \par\smallskip
  {\ERJTableNoteFont 注：$\beta_3$ 表示未加入 $M_{it}$ 时 $Fund_{it}\times Debt$ 对创新的影响；$\theta_1$ 表示 $Fund_{it}\times Debt$ 对 $M_{it}$ 的影响；$\rho_2$ 表示 $M_{it}$ 对创新的影响；$\rho_1$ 为加入 $M_{it}$ 后结果方程中继续保留的 $Fund_{it}\times Debt$ 系数。前两行使用缩尾后标准化的累计基金规模和当期债务压力，后两行使用基金规模对数和滞后一期债务压力对数。括号内为按城市聚类的稳健标准误；$R^2$ 为结果方程拟合优度。***、**、*分别表示在1\%、5\%、10\%水平上显著。}
  }
\end{table}

表8显示，在基准调节项负向显著的前提下，基金规模与债务压力的交互项显著降低存款/GDP口径金融发展水平或其增速，而金融发展变量对创新产出均表现为正向影响并达到10\%显著水平。该结果支持 H5：债务压力会削弱政府产业引导基金改善地区金融发展环境、缓解融资约束的能力，并由此抑制城市创新。需要说明的是，前两行采用缩尾标准化基金规模，后两行采用基金规模对数，因而表8提供的是严格筛选下的机制证据，而非原始基金规模口径的单一检验。

\subsection{本章小结}

本章得到三点结果。第一，政府产业引导基金能够显著提高城市创新产出，基金累计设立规模对发明申请量、实用新型申请量和专利申请总量均为正。第二，地方债务压力显著削弱这一创新促进效应，且当期债务压力和滞后一期债务压力口径下结论一致。第三，机制检验表明，债务压力可能通过降低早期投资事件占比、削弱低金融发展地区的社会资本撬动效率、抑制存款口径金融发展水平三条路径，降低政府产业引导基金的创新绩效。

这些机制结果应按其识别边界理解：耐心资本机制来自基金投资事件数不少于2的分母稳定化样本，社会资本撬动效率机制主要出现在低金融发展地区，融资约束机制属于地区金融发展层面的反向代理证据。总体而言，政府产业引导基金的创新促进作用并非无条件成立，其政策绩效取决于地方债务约束和地区资本市场环境。
% ===== END INPUT: 四_实证结果分析.tex =====



% ===== BEGIN INPUT: 五_结论与政策启示.tex =====
% 本文件由 main.tex 输入，作为“结论与政策启示”章节，不单独编译。

\section{结论与政策启示}

\subsection{主要研究结论}

本文围绕地方债务压力约束下政府产业引导基金与城市创新之间的关系展开分析，重点考察了政府产业引导基金是否能够促进城市创新、地方债务压力是否会削弱这一促进作用，以及这种影响可能通过哪些机制表现出来。结合现有样本和实证结果，本文主要得到以下几方面结论。

第一，政府产业引导基金总体上具有较为明显的创新促进作用。从基准回归结果来看，不论采用基金设立数量还是累计设立规模作为核心解释变量，政府产业引导基金扩张与城市创新产出之间都表现出稳定的正向关系。其中，累计设立规模的结果更加稳健，说明政府产业引导基金并不只是扩大了财政性投入规模，更重要的是它在一定程度上改善了城市创新活动所依赖的资本供给环境，从而提升了创新产出水平。就本文的结果而言，政府产业引导基金在城市层面仍然是一种具有积极作用的创新政策工具。

第二，地方债务压力会显著削弱政府产业引导基金促进创新的边际作用。在将债务压力纳入模型之后，基金累计设立规模与滞后一期债务压力的交互项在发明申请、实用新型申请和专利申请总量等不同口径下均显著为负。这说明，政府产业引导基金的创新促进作用并不是无条件成立的，当地方政府面临较强债务约束时，即便基金规模相近，其推动城市创新的效果也会明显下降。也就是说，地方债务压力构成了政府产业引导基金作用边界中的一个重要约束条件。

第三，债务压力对基金创新效应的削弱，主要可能通过三条路径表现出来。一是耐心资本属性弱化。机制检验表明，在分母稳定化样本中，债务压力会降低基金投向早期项目的比例，而早期投资占比下降又会进一步削弱后续创新产出，这说明债务压力会压缩基金“投早、投小、投长期”的能力。二是社会资本撬动效率下降。现有结果显示，在低金融发展地区，债务压力会削弱社会资本撬动效率，而更高的撬动效率又能够强化基金规模对创新的促进作用，这表明债务约束会在一定程度上削弱财政资金撬动社会资本的政策功能。三是融资约束缓解效应减弱。本文进一步以存款/GDP口径金融发展水平作为融资约束的反向代理，发现债务压力会削弱基金改善地区金融发展环境的能力，而金融发展水平提升又有助于提高专利申请总量对数，因此，基金缓解融资约束的功能在高债务压力情形下也会受到抑制。

第四，本文的机制结论具有一定边界条件，因而更适合被理解为与主结论方向一致的补充证据，而不是对全部情形的无条件推广。比如，耐心资本机制主要建立在基金投资事件数不少于2的样本基础上，社会资本撬动效率机制主要在低金融发展地区更为明显，融资约束缓解机制属于地区金融发展层面的反向代理证据。因此，本文更稳妥的结论是：政府产业引导基金对创新具有正向作用，但其政策绩效在很大程度上取决于地方债务约束和区域资本市场条件，债务压力越高，基金原有的政策性、长期性和引导性功能越容易被压缩。

\subsection{政策启示}

基于上述研究结论，并结合当前关于政府投资基金和创业投资高质量发展的政策导向，本文认为，政府产业引导基金要更好发挥创新促进作用，关键不只是继续扩大基金规模，更在于处理好基金发展、债务风险约束和区域金融环境之间的关系。由此，可以得到以下几方面政策启示。

第一，应进一步突出政府产业引导基金的政策性定位，稳定其耐心资本属性。当前政策文件已经明确提出，政府投资基金要围绕服务国家战略、支持科技创新和发展耐心资本来优化布局和规范运作。结合本文结果看，如果基金运行中过度强调短期收益、安全回收和保值增值，就容易压缩其投向早期项目和高风险创新项目的空间，进而削弱创新促进功能。因此，地方政府在基金考核、出资安排、容错机制和退出节奏设计上，应更加突出长期性和创新导向，避免将政府产业引导基金简单异化为短期财政工具或稳健型投资工具。

第二，应把政府产业引导基金发展与地方债务风险防控统筹起来，避免在高债务压力下片面追求基金扩张。国务院办公厅关于促进政府投资基金高质量发展的指导意见已经强调，基金布局应统筹考虑财力基础、产业基础和债务风险。本文的实证结果进一步说明，债务压力不仅会影响基金设立能力，更会削弱基金运行质量和政策效果。基于此，地方政府在推动产业引导基金设立和扩容时，应更加注重财政承受能力评估和债务风险约束，避免在高债务压力背景下盲目扩张基金规模，防止基金数量上去了，但创新效应反而下降。

第三，应提高基金撬动社会资本的实际效率，减少与市场化资本的同质化竞争。政府产业引导基金的重要政策价值，不仅体现在直接投入，更体现在通过增信、让利和母子基金结构把社会资本引向创新领域。本文结果显示，债务压力会削弱这一撬动功能，尤其在金融环境较弱的地区更加明显。因此，地方政府应进一步完善基金市场化运作机制，强化专业管理人遴选和约束机制，更好发挥财政资金的引导作用，避免基金与市场化资本在中后期成熟项目上重复配置资源，而应更多进入市场资本天然不足的早期科技创新领域。

第四，应将政府产业引导基金与区域金融环境建设结合起来，增强其缓解融资约束的传导基础。本文发现，基金改善存款口径金融发展水平、进而支持创新的作用，在债务约束较强时会被削弱。这说明，仅靠基金本身并不足以完全解决创新融资问题，地方政府还需要将基金发展与科技金融体系建设、直接融资渠道培育、信用环境改善等工作协同推进。特别是在金融发展水平较低、创新融资基础较弱的地区，更应通过基金、银行、担保、资本市场等多种工具的协同配合，增强财政资金支持创新的持续性和传导效率。

第五，应根据地区债务压力和金融发展条件实行差异化基金治理，而不宜采取完全同质化的政策模式。对债务压力较高的地区而言，政策重点应更多放在规范基金运行、强化风险约束和提高资金使用效率上；对金融发展水平较低的地区而言，则应更加重视基金的资本引导、增信补位和社会资本撬动功能；对产业基础较好、创新需求较强的地区，则可以进一步强化基金支持关键核心技术和战略性新兴产业的功能。通过分类施策，有助于避免政府产业引导基金在不同地区出现目标同质化、运作同质化和绩效分化过大的问题。

总体来看，政府产业引导基金仍然是支持区域创新和优化资本配置的重要政策工具，但其创新促进作用并非天然稳定，而是明显受到地方债务压力和区域金融条件的约束。今后推动政府投资基金和创业投资高质量发展，不能只强调基金规模和覆盖面，更需要重视其政策功能的稳定发挥，特别是要在债务风险可控、基金定位清晰和资本市场环境持续改善的条件下，真正把政府产业引导基金建设成为服务长期创新的耐心资本工具。
% ===== END INPUT: 五_结论与政策启示.tex =====


\section*{参考文献}

{\songti\zihao{6}
\ERJRef{蔡庆丰、刘昊、舒少文，2024：《政府产业引导基金与域内企业创新：引导效应还是挤出效应？》，《财贸经济》第3期。}
\ERJRef{陈思霞、陈志勇，2015：《需求回应与地方政府性债务约束机制：经验启示与分析》，《财贸经济》第2期。}
\ERJRef{崔军、花培严、王泓澄，2025：《政府引导基金扩张之谜》，《厦门大学学报（哲学社会科学版）》第3期。}
\ERJRef{崔竹，2025：《地方政府投资基金发展现状、风险特征及监管对策》，《行政管理改革》第4期。}
\ERJRef{国务院，2025：《关于2024年度政府债务管理情况的报告》，中国人大网。}
\ERJRef{国务院办公厅，2024：《国务院办公厅关于印发〈促进创业投资高质量发展的若干政策措施〉的通知》，中国政府网，2024年6月19日，\url{https://www.gov.cn/zhengce/zhengceku/202406/content_6958231.htm}。}
\ERJRef{国务院办公厅，2025：《国务院办公厅关于促进政府投资基金高质量发展的指导意见》，中国政府网，2025年1月7日，\url{https://www.gov.cn/zhengce/zhengceku/202501/content_6996730.htm}。}
\ERJRef{华夏、马树才、韩云虹，2020：《地方政府债务如何影响实体企业信贷融资——基于异质性视角的中国工业企业微观数据分析》，《贵州财经大学学报》第3期。}
\ERJRef{孔丹凤、谢国梁，2021：《地方政府债务、债务结构与企业金融化》，《现代经济探讨》第11期。}
\ERJRef{李健旋、赵林度，2018：《金融集聚、生产率增长与城乡收入差距的实证分析——基于动态空间面板模型》，《中国管理科学》第12期。}
\ERJRef{李森、王聪，2024：《财政压力对制造业企业创新的影响研究——基于创新数量与创新质量双重视角的经验分析》，《财经科学》第4期。}
\ERJRef{刘贯春、程飞阳、姚守宇、张军，2022：《地方政府债务治理与企业投融资期限错配改善》，《管理世界》第11期。}
\ERJRef{刘玉斌、赵天宇、郭树龙，2023：《战略性新兴产业创业投资引导基金能促进企业创新吗？》，《产业经济研究》第1期。}
\ERJRef{罗志恒，2021：《政府引导基金：发展历程、运行机制及对隐性债务的影响》，粤开证券研究报告。}
\ERJRef{施国平、党兴华、董建卫，2016：《引导基金能引导创投机构投向早期和高科技企业吗？——基于双重差分模型的实证评估》，《科学学研究》第6期。}
\ERJRef{王平、徐肇仪，2023：《数字普惠金融、融资约束与共同富裕》，《会计之友》第10期。}
\ERJRef{吴超鹏、严泽浩，2023：《政府基金引导与企业核心技术突破：机制与效应》，《经济研究》第6期。}
\ERJRef{徐明，2022：《政府风险投资、代理问题与企业创新——来自政府引导基金介入的证据》，《南开经济研究》第2期。}
\ERJRef{徐彦坤，2020：《地方政府债务如何影响企业投融资行为？》，《中南财经政法大学学报》第2期。}
\ERJRef{余海跃、康书隆，2020：《地方政府债务扩张、企业融资成本与投资挤出效应》，《世界经济》第7期。}
\ERJRef{张杰，2020：《政府创新补贴对中国企业创新的激励效应——基于U型关系的一个解释》，《经济研究》第6期。}
\ERJRef{张杰、任元明，2025：《探寻以耐心资本为内核的中国特色现代金融体系发展路径与改革突破口》，《天津社会科学》第2期。}
\ERJRef{张壹帆、陆岷峰，2025：《新质生产力视角下政府引导基金“耐心度”培育：评价体系与利益平衡》，《证券市场导报》第1期。}
\ERJRef{张跃文、李慧、徐枫，2025：《财政压力、“合作型”政府干预与企业投资》，《财经论丛》第9期。}
\ERJRef{张圆圆、夏球，2024：《我国政府引导基金创新发展研究》，《西南金融》第3期。}
\ERJRef{周天宇，2024：《地方政府债务变动对企业创新的影响研究——基于333个地级行政区上市企业的数据》，《中国市场》第10期。}
\ERJRef{周伟、陆佳玮，2022：《财政压力、地方政府行为与企业创新》，出版信息待核。}
\ERJRef{Cumming, D. J., 2007, “Government Policy towards Entrepreneurial Finance: Innovation Investment Funds”, Journal of Business Venturing, Vol. 22, No. 2, pp. 193-235.}
\ERJRef{Cumming, D. J., and S. Johan, 2007, “Pre-Seed Government Venture Capital Funds”, Journal of International Entrepreneurship, forthcoming.}
\ERJRef{Fazzari, S., and B. Petersen, 1993, “Working Capital and Fixed Investment: New Evidence on Financing Constraints”, RAND Journal of Economics, Vol. 24, pp. 328-342.}
\ERJRef{Hall, B., 2002, “The Financing of Research and Development”, Oxford Review of Economic Policy, Vol. 18, pp. 35-51.}
\ERJRef{Hall, B., 2005, “The Financing of Innovation”, in S. Shane (ed.), Blackwell Handbook of Technology and Innovation Management, Blackwell Publishers, Oxford.}
\ERJRef{Hao, P., Y. Wang, and L. Fan, 2025, “Government-guided Fund, Social Resources, and Corporate Green Innovation”, International Review of Financial Analysis, Vol. 98.}
\ERJRef{Howell, S. T., 2024, “Government Intervention in Innovation”, Annual Review of Financial Economics, Vol. 16, pp. 367-390.}
\ERJRef{Manso, G., 2011, “Motivating Innovation”, The Journal of Finance, Vol. 66, No. 5, pp. 1823-1860.}
\ERJRef{Wang, Y., et al., 2025, “Does Local Government Debt Affect Corporate Innovation Quality? Evidence from China”, Sustainability, Vol. 17, No. 2, p. 550.}
\ERJRef{Xiao, F., Z. Liang, Y. Lv, et al., 2024, “The Effect of Government-guided Funds on Target Industries in Development Zones: Evidence from China”, Accounting and Finance, Vol. 64, No. 5, pp. 4701-4722.}
\ERJRef{Zhang, X., and R. Jin, 2022, “Has Local Government Debt Crowded Out Enterprise Innovation?”, PLOS ONE, Vol. 17, No. 11, e0277461.}
\ERJRef{Zhao, C., and L. Zhang, 2025, “Can Government-guided Funds Promote Innovation Output in Strategic Emerging Enterprises? Evidence from China”, PLOS ONE, Vol. 20, No. 10, e0334826.}
}

\vspace{2\baselineskip}


\end{document}
